\documentclass[11pt,a4paper]{article}
\usepackage[utf8]{inputenc}
\usepackage[T1]{fontenc}
\usepackage{lmodern}
\usepackage{microtype}
\usepackage[margin=2.2cm]{geometry}
\usepackage{enumitem}
\setlist{nosep,leftmargin=1.5em}
\setlength{\parindent}{0pt}
\setlength{\parskip}{0.5em}
\pagestyle{empty}

\newcommand{\polje}[1]{\textbf{#1}}

\begin{document}

{\large\textbf{Prijava teme seminarske naloge}}\\[0.2em]
Nekonvencionalne platforme in metode procesiranja, 2026/2027\\
Skupina 08: Jurij Anžič, Janez Koprivec

\vspace{0.6em}

\polje{Naslov:} Primerjava kvantnih seštevalnikov CDKM in Draper (QFT) na šumnem
kvantnem procesorju IBM

\polje{Opis:} V Qiskitu bomo implementirala Cuccarov seštevalnik z zaporednim
prenosom (CDKM) in Draperjev seštevalnik na osnovi kvantne Fourierjeve
transformacije (QFT). Oba bosta seštevala dve $n$-bitni števili v
$(n{+}1)$-bitni rezultat. Za $n = 1$ do $5$ ju bomo izvedla na idealnem
simulatorju, na simulatorju s šumnim modelom in na procesorju IBM.

\polje{Raziskovalno vprašanje:} Kateri seštevalnik na procesorju IBM pogosteje
vrne pravilno vsoto in ali to napove število dvokubitnih vrat po prevajanju?

\polje{Hipoteza:} Pri vsakem $n$ od $2$ do $5$ bo CDKM na procesorju IBM imel
višjo verjetnost pravilne vsote kot Draper, ker bo imel po prevajanju manj
dvokubitnih vrat.

\polje{Metrike:}
\begin{itemize}
  \item verjetnost pravilne vsote, primerjana z naključnim izidom ($1/2^{n+1}$)
  \item število kubitov, globina in število dvokubitnih vrat pred prevajanjem in po njem
\end{itemize}

\polje{Načrt poskusov:}
\begin{enumerate}
  \item Preverjanje pravilnosti obeh seštevalnikov na idealnem simulatorju.
  \item Prevajanje za izbrani procesor IBM z enakimi nastavitvami za oba in
        zapis lastnosti vezij.
  \item Zagon na šumnem simulatorju in nato na procesorju IBM z enakimi pari
        vhodov za oba seštevalnika (vsi pari pri $n \le 2$, sicer 16 naključnih).
  \item Primerjava uspešnosti in števila vrat.
\end{enumerate}

\polje{Orodja:} Qiskit, Qiskit Aer, IBM Quantum (Open Plan), matplotlib,
Overleaf, GitHub.

\polje{Razdelitev dela:}
\begin{itemize}
  \item \emph{Posebej:} vsak od članov posebej pregleda enega od seštevalnikov
        in literaturo zanj.
  \item \emph{Skupaj:} zagon poskusov, primerjava in grafi, poročilo in predstavitev.
\end{itemize}

\end{document}
